\section{统一建模框架}

\subsection{基本符号}

\begin{table}[H]
\centering
\caption{核心符号与含义}
\label{tab:symbols}
\begin{tabularx}{0.96\textwidth}{C{1.8cm}Y C{1.8cm}Y}
\toprule
符号 & 含义 & 符号 & 含义\\
\midrule
$G=(V,E)$ & 区域邻接图 & $T$ & 截止日期\\
$z_t$ & 第 $t$ 天天气 & $l_t$ & 日末位置\\
$w_t,f_t$ & 日末水、食物库存 & $m_t$ & 日末现金\\
$a_t$ & 移动、停留或挖矿行动 & $b^w_t,b^f_t$ & 当日购买量\\
$c^r_z$ & 资源 $r$ 在天气 $z$ 下的基础消耗 & $\lambda_t$ & 行动消耗倍数\\
$B$ & 矿山基础收益 & $C$ & 负重上限\\
$n,k$ & 总玩家数、同组行动人数 & $J_i$ & 玩家 $i$ 的终值\\
\bottomrule
\end{tabularx}
\end{table}

\subsection{状态转移和收益}

单人状态写为 $s_t=(t,l_t,w_t,f_t,m_t)$。设 $r\in\{w,f\}$，行动倍数为
\[
\lambda(a_t)=
\begin{cases}
2,&a_t=\text{移动},\\
1,&a_t=\text{停留},\\
3,&a_t=\text{挖矿}.
\end{cases}
\]
则不含购买时的库存转移为
\begin{equation}
r_t=r_{t-1}-\lambda(a_t)c^r_{z_t},\qquad r_t\ge0,
\label{eq:resource}
\end{equation}
并在每日采购后满足
\begin{equation}
3w_t+2f_t\le1200.
\label{eq:capacity}
\end{equation}
若在矿山挖矿，现金增加 $B$；村庄购买单价是基准价两倍；终点终值为
\begin{equation}
J=m_\tau+\tfrac12p_ww_\tau+\tfrac12p_ff_\tau,
\label{eq:terminal}
\end{equation}
其中 $\tau\le T$ 为首次到达日。

\subsection{证据等级与统一验收}

为避免把不同强度的计算结论混在一起，本文采用四级证据：

\begin{enumerate}
  \item \evidence{E1} 求解器最优证书与独立回放，对应“全局最优”；
  \item \evidence{E2} 完整状态或天气支撑精确枚举，对应“给定模型下精确最优”；
  \item \evidence{E3} 完整最佳响应检验，对应“纯策略均衡”；
  \item \evidence{E4} 独立种子样本外评估与区间，对应“条件推荐”。
\end{enumerate}

所有策略都经过独立规则入口回放。Monte Carlo 中失败样本的终值统一记为0，成功样本均值只用于解释“到达后的财富”，选模和结论使用含失败的总体效用。

\begin{resultbox}{统一边界}
问题一可以证明全局最优；第三、五关可以在给定有限支撑或策略空间内精确验证；第四、六关因天气分布未给，只能在声明分布和风险预算下推荐。四类表述不可互换。
\end{resultbox}
